\documentclass[%
    10pt,
    french,
    aspectratio=1610,
    xcolor=svgnames
]{beamer}

\usetheme[
    decorative,
    nocodebox,
    frametitle=plain,
    margin=3em,
    soberblock
]{Celestia}

\title{Le thème Celestia}
\subtitle{v1.0.0}
\author{Razik Ikhlef}
\date{\today}

\begin{document}

\begin{frame}[plain]
    \titlepage
\end{frame}

\begin{frame}[fragile]{Polices}
    \begin{description}[mainfaceoptions]
        \item[mainface] Police principale pour le texte et les titres (\alert{Literata} par défaut)
        \item[mainfaceoptions] Options directement passées à \alert{fontspec} pour la police principale
        \item[sansface] Police sans empattements pour les éléments structurels (\alert{Inter} par défaut)
        \item[sansfaceoptions] Options passées à \alert{fontspec} pour la police sans empattements
        \item[monoface] Police à chasse fixe pour le code (\alert{Roboto Mono} par défaut)
        \item[monofaceoptions] Options passées à \alert{fontspec} pour la police à chasse fixe
        \item[allserif] Utilise la police principale avec empattements pour les mathématiques
    \end{description}

    \begin{exampleblock}{Exemple}
        \begin{lstlisting}[style=latex]
\usetheme[
    mainface=EB Garamond,
    mainfaceoptions={Scale=1.1},
    sansface=Montserrat,
    monoface=Fira Code,
    allserif
]{Celestia}
        \end{lstlisting}
    \end{exampleblock}
\end{frame}

\begin{frame}[fragile]{Couleurs}
    \begin{description}[codebackgroundcolor]
        \item[maincolor] Couleur principale utilisée pour les titres et accents (code \alert{\LaTeX} svgname ou \alert{HTML} hexadécimal)
        \item[accentcolor] Couleur secondaire pour les éléments spéciaux
        \item[backgroundcolor] Couleur de fond des diapositives (\alert{F7F9FC} par défaut)
        \item[codebackgroundcolor] Couleur de fond des blocs de code (\alert{F1F3F6} par défaut)
        \item[mainblue] Couleur des blocs standards (\alert{045549} par défaut)
        \item[maingreen] Couleur des blocs exemple (\alert{054924} par défaut)
        \item[mainred] Couleur des blocs alerte (\alert{490445} par défaut)
        \item[unicolor] Utilise la couleur principale pour tout le texte
    \end{description}

    \begin{exampleblock}{Exemple}
        \begin{lstlisting}[style=latex]
\usetheme[
    maincolor=045549,
    accentcolor=E63946,
    backgroundcolor=FAFAFA,
    codebackgroundcolor=F5F5F5
]{Celestia}
        \end{lstlisting}
    \end{exampleblock}
\end{frame}

\begin{frame}[fragile]{Mise en page}
    \begin{description}[centeredtitle]
        \item[margin] Marge du contenu (\alert{2em} par défaut)
        \item[frametitle] Style du titre (\alert{elegant}, \alert{plain}, \alert{centered})
        \item[decorative] Active les éléments décoratifs
        \item[nodecorative] Désactive les éléments décoratifs
        \item[centeredtitle] Centre le titre sur la page de titre
        \item[titleright] Aligne le titre à droite
    \end{description}

    \begin{exampleblock}{Exemple}
        \begin{lstlisting}[style=latex]
\usetheme[
    margin=1.5em,
    frametitle=centered,
    decorative,
    centeredtitle
]{Celestia}
        \end{lstlisting}
    \end{exampleblock}
\end{frame}

\begin{frame}[fragile]{En-têtes}
    \begin{description}[headweight]
        \item[headstyle] Famille de police pour les titres : \alert{rmfamily} (avec empattements) ou \alert{sffamily} (sans empattements)
        \item[headshape] Style des caractères : \alert{sc} (petites capitales), \alert{it} (italique), \alert{normal}
        \item[headweight] Graisse des titres : \alert{bfseries} (gras) ou \alert{mdseries} (normal)
    \end{description}

    \begin{exampleblock}{Exemple}
        \begin{lstlisting}[style=latex]
\usetheme[
    headstyle=sffamily,
    headshape=sc,
    headweight=bfseries
]{Celestia}
        \end{lstlisting}
    \end{exampleblock}
\end{frame}

\begin{frame}[fragile]{Code}
    \begin{description}[nocodeframe]
        \item[codebox] Active l'encadrement du code avec \alert{tcolorbox} (\alert{true} par défaut)
        \item[nocodebox] Désactive complètement l'encadrement \alert{tcolorbox} du code
        \item[nocodeframe] Conserve \alert{tcolorbox} mais sans bordure visible
    \end{description}

    \begin{exampleblock}{Exemple}
        \begin{lstlisting}[style=latex]
\usetheme[
    nocodebox,
    nocodeframe
]{Celestia}
        \end{lstlisting}
    \end{exampleblock}
\end{frame}

\begin{frame}[fragile]{Pied de page}
    \begin{description}[quartercirclefooter]
        \item[nofooter] Supprime entièrement le pied de page, sauf le numéro de diapositive
        \item[quartercirclefooter] Affiche uniquement le numéro dans un quart de cercle en bas à droite
        \item[fullbarfooter] Crée une barre complète avec auteur/titre/date et numéro dans un cercle
    \end{description}

    \begin{exampleblock}{Exemple}
        \begin{lstlisting}[style=latex]
\usetheme[
    quartercirclefooter
]{Celestia}
        \end{lstlisting}
    \end{exampleblock}
\end{frame}

\begin{frame}[fragile]{Table des matières}
    \begin{description}[twocolumntoc]
        \item[compacttoc] Réduit l'espacement vertical entre les entrées de la table
        \item[twocolumntoc] Répartit automatiquement les sections sur deux colonnes équilibrées
    \end{description}

    \begin{exampleblock}{Exemple}
        \begin{lstlisting}[style=latex]
\usetheme[
    compacttoc,
    twocolumntoc
]{Celestia}
        \end{lstlisting}
    \end{exampleblock}
\end{frame}

\begin{frame}{Blocs}
    \begin{description}[exampleblock]
        \item[block] Bloc standard pour le contenu normal
        \item[exampleblock] Bloc pour les exemples
        \item[alertblock] Bloc pour les alertes
    \end{description}

    \begin{block}{Bloc standard}Contenu d'un bloc standard\end{block}
    \begin{exampleblock}{Bloc exemple}Contenu d'un bloc exemple\end{exampleblock}
    \begin{alertblock}{Bloc alerte}Contenu d'un bloc alerte\end{alertblock}
\end{frame}

\begin{frame}[fragile]{Style des blocs}
    \begin{description}[soberblock]
        \item[soberblock] Le titre adopte la couleur principale du bloc (mainblue/maingreen/mainred) sur le fond général du document, tandis que le corps garde un fond légèrement teinté (10\%)
        \item[softblock] Le titre et le corps partagent le même fond légèrement teinté (10\%), avec le titre dans la couleur principale correspondante
    \end{description}

    \begin{exampleblock}{Exemple}
        \begin{lstlisting}[style=latex]
\usetheme[
    soberblock  % Titre coloré sur fond document
]{Celestia}

% ou

\usetheme[
    softblock   % Titre coloré sur fond 10%
]{Celestia}
        \end{lstlisting}
    \end{exampleblock}

\end{frame}

\begin{frame}[standout]{Pages d'emphase}
    L'option \emph{standout} transforme une diapositive en page d'emphase, idéale pour les moments clés de la présentation : citations marquantes, chiffres essentiels, messages à retenir
\end{frame}

\end{document}
